\chapter{2016年天津卷（理）}

\section{单选题}

\begin{problem}
已知集合 \(A=\{1,2,3,4\}\)，\(B=\{y\mid y=3x-2,\ x\in A\}\)，则 \(A\cap B=\)（\quad）

\choices
    {\(\{1\}\)}
    {\(\{4\}\)}
    {\(\{1,3\}\)}
    {\(\{1,4\}\)}
\end{problem}

\begin{problem}
设变量 \(x\)，\(y\) 满足约束条件
\(\begin{cases}
x-y+2\ge 0,\\
2x+3y-6\ge 0,\\
3x+2y-9\le 0,
\end{cases}\)
则目标函数 \(z=2x+5y\) 的最小值为（\quad）

\choices
    {\(-4\)}
    {\(6\)}
    {\(10\)}
    {\(17\)}
\end{problem}

\begin{problem}
在 \(\triangle ABC\) 中，若 \(AB=\sqrt{13}\)，\(BC=3\)，\(\angle C=120^\circ\)，则 \(AC=\)（\quad）

\choices
    {\(1\)}
    {\(2\)}
    {\(3\)}
    {\(4\)}
\end{problem}

\begin{problem}
阅读如下的程序框图，运行相应的程序，则输出 \(S\) 的值为（\quad）

\texfigure[max width=0.38\linewidth]{img_repaint/2016/tianjin_science/q4_fig1.tex}

\choices
    {\(2\)}
    {\(4\)}
    {\(6\)}
    {\(8\)}
\end{problem}

\begin{problem}
设 \(\{a_n\}\) 是首项为正数的等比数列，公比为 \(q\)，则“\(q<0\)”是“对任意的正整数 \(n\)，\(a_{2n-1}+a_{2n}<0\)”的（\quad）

\choices
    {充要条件}
    {充分而不必要条件}
    {必要而不充分条件}
    {既不充分也不必要条件}
\end{problem}

\begin{problem}
已知双曲线 \(\dfrac{x^2}{4}-\dfrac{y^2}{b^2}=1\)（\(b>0\)），以原点为圆心，双曲线的实半轴长为半径长的圆与双曲线的两条渐近线相交于 \(A\)，\(B\)，\(C\)，\(D\) 四点，四边形 \(ABCD\) 的面积为 \(2b\)，则双曲线的方程为（\quad）

\choices
    {\(\dfrac{x^2}{4}-\dfrac{3y^2}{4}=1\)}
    {\(\dfrac{x^2}{4}-\dfrac{4y^2}{3}=1\)}
    {\(\dfrac{x^2}{4}-\dfrac{y^2}{4}=1\)}
    {\(\dfrac{x^2}{4}-\dfrac{y^2}{12}=1\)}
\end{problem}

\begin{problem}
已知 \(\triangle ABC\) 是边长为 \(1\) 的等边三角形，点 \(D\)，\(E\) 分别是边 \(AB\)，\(BC\) 的中点，连接 \(DE\) 并延长到点 \(F\)，使得 \(DE=2EF\)，则 \(\overrightarrow{AF}\cdot\overrightarrow{BC}\) 的值为（\quad）

\choices
    {\(-\dfrac85\)}
    {\(\dfrac18\)}
    {\(\dfrac14\)}
    {\(\dfrac{11}{8}\)}
\end{problem}

\begin{problem}
已知函数 \(f(x)=\begin{cases}
x^2+(4a-3)x+3a,&x<0,\\
\log_a(x+1)+1,&x\ge 0,
\end{cases}\)（\(a>0\)，且 \(a\neq1\)）在 \(\R\) 上单调递减，且关于 \(x\) 的方程 \(\lvert f(x)\rvert=2-x\) 恰好有两个不相等的实数解，则 \(a\) 的取值范围是（\quad）

\choices
    {\(\left(0,\dfrac23\right]\)}
    {\(\left[\dfrac23,\dfrac34\right]\)}
    {\(\left[\dfrac13,\dfrac23\right]\cup\left\{\dfrac34\right\}\)}
    {\(\left[\dfrac13,\dfrac23\right)\cup\left\{\dfrac34\right\}\)}
\end{problem}

\section{填空题}

\begin{problem}
已知 \(a\)，\(b\in\R\)，\(i\) 是虚数单位，若 \((1+i)(1-bi)=a\)，则 \(\dfrac{a}{b}\) 的值为\fillinblank{}.
\end{problem}

\begin{problem}
\(\left(x^2-\dfrac1x\right)^8\) 的展开式中 \(x^7\) 的系数为\fillinblank{}（用数字作答）.
\end{problem}

\begin{problem}
已知一个四棱锥的底面是平行四边形，该四棱锥的三视图如图所示（单位：\(\mathrm{m}\)），则该四棱锥的体积为\fillinblank{}\(\mathrm{m}^3\).

\bitmapfigure[max width=0.58\linewidth]{img/2016/tianjin_science/q11_fig1.png}
\end{problem}

\begin{problem}
如图，\(AB\) 是圆的直径，弦 \(CD\) 与 \(AB\) 相交于点 \(E\)，\(BE=2AE=2\)，\(BD=ED\)，则线段 \(CE\) 的长为\fillinblank{}.

\bitmapfigure[max width=0.34\linewidth]{img/2016/tianjin_science/q12_fig1.png}
\end{problem}

\begin{problem}
已知 \(f(x)\) 是定义在 \(\R\) 上的偶函数，且在区间 \((-\infty,0)\) 上单调递增. 若实数 \(a\) 满足 \(f\left(2^{\lvert a-1\rvert}\right)>f\left(-\sqrt2\right)\)，则 \(a\) 的取值范围是\fillinblank{}.
\end{problem}

\begin{problem}
设抛物线
\(\begin{cases}
x=2pt^2,\\
y=2pt,
\end{cases}\quad(t\text{ 为参数},\ p>0)\)
的焦点为 \(F\)，准线为 \(l\). 过抛物线上一点 \(A\) 作 \(l\) 的垂线，垂足为 \(B\). 设 \(C\left(\dfrac72p,0\right)\)，\(AF\) 与 \(BC\) 相交于点 \(E\). 若 \(\lvert CF\rvert=2\lvert AF\rvert\)，且 \(\triangle ACE\) 的面积为 \(3\sqrt2\)，则 \(p\) 的值为\fillinblank{}.
\end{problem}

\section{解答题}

\begin{problem}
已知函数 \(f(x)=4\tan x\sin\left(\dfrac\pi2-x\right)\cos\left(x-\dfrac\pi3\right)-\sqrt3\).

\begin{enumerate}
\item 求 \(f(x)\) 的定义域与最小正周期；
\item 讨论 \(f(x)\) 在区间 \(\left[-\dfrac\pi4,\dfrac\pi4\right]\) 上的单调性.
\end{enumerate}
\end{problem}

\begin{problem}
某小组共 \(10\) 人，利用假期参加义工活动，已知参加义工活动次数为 \(1\)，\(2\)，\(3\) 的人数分别为 \(3\)，\(3\)，\(4\). 现从这 \(10\) 人中随机选出 \(2\) 人作为该组代表参加座谈会.

\begin{enumerate}
\item 设 \(A\) 为事件“选出的 \(2\) 人参加义工活动次数之和为 \(4\)”，求事件 \(A\) 发生的概率；
\item 设 \(X\) 为选出的 \(2\) 人参加义工活动次数之差的绝对值，求随机变量 \(X\) 的分布列和数学期望.
\end{enumerate}
\end{problem}

\begin{problem}
如图，正方形 \(ABCD\) 的中心为 \(O\)，四边形 \(OBEF\) 为矩形，平面 \(OBEF\perp\) 平面 \(ABCD\)，点 \(G\) 为 \(AB\) 的中点，\(AB=BE=2\).

\begin{enumerate}
\item 求证：\(EG\parallel\) 平面 \(ADF\)；
\item 求二面角 \(O-EF-C\) 的正弦值；
\item 设 \(H\) 为线段 \(AF\) 上的点，且 \(AH=\dfrac23HF\)，求直线 \(BH\) 和平面 \(CEF\) 所成角的正弦值.
\end{enumerate}

\bitmapfigure[max width=0.55\linewidth]{img/2016/tianjin_science/q17_fig1.png}
\end{problem}

\begin{problem}
已知 \(\{a_n\}\) 是各项均为正数的等差数列，公差为 \(d\)，对任意的 \(n\in\N^*\)，\(b_n\) 是 \(a_n\) 和 \(a_{n+1}\) 的等比中项.

\begin{enumerate}
\item 设 \(c_n=b_{n+1}^2-b_n^2\)（\(n\in\N^*\)），求证：数列 \(\{c_n\}\) 是等差数列；
\item 设 \(a_1=d\)，\(T_n=\displaystyle\sum_{k=1}^{2n}(-1)^k b_k^2\)（\(n\in\N^*\)），求证：\(\displaystyle\sum_{k=1}^n\dfrac1{T_k}<\dfrac1{2d^2}\).
\end{enumerate}
\end{problem}

\begin{problem}
设椭圆 \(\dfrac{x^2}{a^2}+\dfrac{y^2}{3}=1\)（\(a>\sqrt3\)）的右焦点为 \(F\)，右顶点为 \(A\)，已知 \(\dfrac1{\lvert OF\rvert}+\dfrac1{\lvert OA\rvert}=\dfrac{3e}{\lvert FA\rvert}\)，其中 \(O\) 为原点，\(e\) 为椭圆的离心率.

\begin{enumerate}
\item 求椭圆的方程；
\item 设过点 \(A\) 的直线 \(l\) 与椭圆交于点 \(B\)（\(B\) 不在 \(x\) 轴上），垂直于 \(l\) 的直线与 \(l\) 交于点 \(M\)，与 \(y\) 轴交于点 \(H\). 若 \(BF\perp HF\)，且 \(\angle MOA\le\angle MAO\)，求直线 \(l\) 的斜率的取值范围.
\end{enumerate}
\end{problem}

\begin{problem}
设函数 \(f(x)=(x-1)^3-ax-b\)（\(x\in\R\)，其中 \(a\)，\(b\in\R\)）.

\begin{enumerate}
\item 求 \(f(x)\) 的单调区间；
\item 若 \(f(x)\) 存在极值点 \(x_0\)，且 \(f(x_1)=f(x_0)\)，其中 \(x_1\neq x_0\)，求证：\(x_1+2x_0=3\)；
\item 设 \(a>0\)，函数 \(g(x)=\lvert f(x)\rvert\)，求证：\(g(x)\) 在区间 \([0,2]\) 上的最大值不小于 \(\dfrac14\).
\end{enumerate}
\end{problem}
